\documentclass[10pt]{article}
\usepackage[margin=20mm]{geometry}
\usepackage{amsmath,amssymb,mathtools,fontspec,unicode-math,xurl,hyperref}
\IfFontExistsTF{DejaVu Serif}{\setmainfont{DejaVu Serif}}{\setmainfont{Latin Modern Roman}}
\setmathfont{latinmodern-math.otf}
\setlength{\parindent}{0pt}\setlength{\parskip}{6pt}\sloppy
\title{Technical Review of
Renormalization Group and Elliptic Homogenization in High Contrast}
\author{AI-assisted technical review}
\date{}
\begin{document}
\maketitle
Abstract: The reviewed source is the 152-page author manuscript arXiv:2405.10732v3. Seven substantive concerns are established within the inspected scope. They include an admissible counterexample to the full Sobolev-exponent claim of Theorem B, a sharp matrix counterexample to Lemma 2.6, failures of two advertised coefficient-field examples, an incorrect independent-sum tail bound, and two invalid proof dependencies. The report supplies explicit assumptions, calculations, counterevidence, and bounded repairs where proved. These findings support Major revision of the reviewed version; they do not establish that the uniformly elliptic homogenization theorem or the complete renormalization program is false. Scientific adjudication remains pending.\par
\section*{Identity, version, and review boundary}
The paper is by Scott Armstrong and Tuomo Kuusi. The scientific source for every numbered location below is arXiv:2405.10732v3, 152 PDF pages, bearing the internal date September 5, 2025 and the arXiv watermark September 6, 2025 [AKv3]. Printed and PDF page numbers agree. The publisher record identifies Inventiones mathematicae 242 (2025), 895–1086, DOI https://doi.org/10.1007/s00222-025-01370-9. The publisher’s full journal PDF was not obtained: the accessible record is a subscription preview. No mathematical or typographic equivalence between the reviewed author manuscript and the journal version is asserted.\par
The prior preserved index contains only source, identity, and manifest materials for this paper, with no recoverable substantive report. The current review therefore supplies the necessary initial examination. It is not independent acceptance and does not adjudicate the paper on behalf of a human editor. Its mathematical objections are established by analytic arguments or exact reviewer-authored arithmetic. No author code was run. Lack of optional code is not a mathematical objection.\par
\section*{Assumptions that govern the checks}
The coefficient field is allowed to be nonsymmetric and microscopically degenerate. Its symmetric and skew parts are \(\mathbf s=(\mathbf a+\mathbf a^t)/2\) and \(\mathbf k=\mathbf a-\mathbf s\). The ambient space on p. 7 requires \[\mathbf s,\quad \mathbf s^{-1},\quad
 \mathbf k^t\mathbf s^{-1}\mathbf k\in L^1_{\rm loc},\] with positive symmetric part, up to irrelevant null-set representatives. Theorem B assumes \(\mathbb Z^d\) stationarity (P1), coarse ellipticity (P2) with \(\gamma\in[0,1)\), and the concentration condition (P3) with \(\nu>\gamma\). The sum in (1.17) is an average; the convention is explicit on p. 15. Treating it as an unnormalized sum would create a spurious objection.\par
The weighted solution spaces \(H^1_{\mathbf a}\) in (2.2)–(2.5) need not embed in ordinary \(H^1\). In contrast, the \(H^{-s}\) norms in (1.20) are ordinary fractional Sobolev dual norms, volume-normalized as in (1.32)–(1.33). This distinction is essential in Problem 1. For the matrix checks, coarse \(\mathbf k(U)\) is not required to be skew-symmetric; this is stated on p. 17 and allowed in Lemma 2.6.\par
\section*{A principal conclusion with an inadmissibly broad exponent range}
Problem 1 (P01: Theorem B, Dirichlet estimate (1.20); Major).  On p. 10, Theorem B asserts its gradient and flux error estimate for every \(s\in(0,\tfrac12)\). That quantifier is false under (P1)–(P3).\par
\subsection*{An admissible field}
Take \(d=2\), \(\alpha=4/5\), \(\gamma=9/10\), \(\nu=1\), and \(\beta=0\). Define the one-periodic function \[w(t)=\bigl(2\operatorname{dist}(t,\mathbb Z)\bigr)^\alpha\quad(t\notin\mathbb Z),
 \qquad w(t)=1\quad(t\in\mathbb Z),
 \qquad
 \mathbf a(x)=\operatorname{diag}(w(x_1),1).\] Changing the values on the integer planes changes no integral or weak equation. The field is bounded, symmetric, positive definite with this representative, and its inverse is locally integrable because \(\alpha<1\). It belongs to \(\Omega\) and has \(\mathbf k=0\), so \(H^1_{\mathbf a}=H^1_{\mathbf s}\).\par
Let the law be the Dirac measure at this field. It is \(\mathbb Z^2\)-stationary. Every centered random variable is zero almost surely, so (P3) holds with the parameters above and \(\Psi(t)=\exp(t^2)\), \(K_\Psi=3\). The additional measurability and sensitivity requirements only restrict which centered variables are tested; they do not affect this observation.\par
For every axis-aligned cube \(Q\) of side length \(r\), the variational bound (2.21), p. 22, implies \[\mathbf b(Q)\le \frac1{|Q|}\int_Q\mathbf a\le I,\qquad
 \mathbf s_*(Q)^{-1}\le\frac1{|Q|}\int_Q\mathbf a^{-1}
 \le 30\max(1,r^{-\alpha})I.\] The scalar estimate follows by integrating the singularity \(|t|^{-\alpha}\) over intervals of length at most one and by covering longer intervals by periods; the one-period average of \(w^{-1}\) is \(5\). In this symmetric case the coarse skew block vanishes. For \(m\ge0\) and \(k\le m\), the inverse-ellipticity term in (1.11) is bounded by \[30\,3^{-\gamma(m-k)}\max(1,3^{-\alpha k})\le30.\] For \(k<0\) this uses \(3^{-\gamma m}3^{(\gamma-\alpha)k}\le1\); for \(k\ge0\) it is immediate. Thus (P2) holds with \(\Lambda_0=1\), \(\lambda_0=1/30\), and \(\mathcal S=1\). One may take \(\Psi_{\mathcal S}(t)=\exp((t-1)_+^2)\) and \(K_{\Psi_{\mathcal S}}=3\): the required tail bound is immediate, and \(\log t+(t-1)^2\le(3t-1)^2\) for \(t\ge1\) proves the growth condition. In particular, \(\nu>\gamma\) is satisfied.\par
\subsection*{A smooth datum with a singular weighted solution}
Set \(L=1/4\), \(U=(-L,L)^2\), \(q=\pi/(2L)\), \(b(x_2)=\cos(qx_2)\), and \(g(x)=x_1b(x_2)\). The function \(g\) is smooth on \(\overline U\). Every ellipsoid \(E_1\) in Theorem A contains the unit Euclidean ball, since \(\overline{\mathbf s}\ge\bar\lambda I\); therefore \(U\subset E_1\).\par
Equation reference: eq:volterra. Fix any \(\varepsilon>0\) and put \(w_\varepsilon(t)=w(t/\varepsilon)\). On \([0,L]\) solve \[
 v(t)=\int_0^t\frac{F(r)}{w_\varepsilon(r)}\,dr,\qquad
 F(t)=c+q^2\int_0^t v(r)\,dr.\] The coefficient \(w_\varepsilon^{-1}\) is integrable, so this linear Volterra system has a unique absolutely continuous solution. For \(c=1\), \(v(L)>0\). By linearity choose \(c=L/v_{c=1}(L)>0\), so that \(v(L)=L\). Extend \(v\) oddly and \(F\) evenly. Then \[-(w_\varepsilon v')'+q^2v=0,\qquad v(\pm L)=\pm L,
 \qquad u^\varepsilon(x)=v(x_1)b(x_2)\] solves the required equation and has boundary data \(g\).\par
The weighted energy is finite because \(\int w_\varepsilon(v')^2=\int F^2/w_\varepsilon<\infty\). Here is also a direct check of the precise zero-boundary closure. Write \(z=v-t\). Approximate \(z'\) in \(L^2(w_\varepsilon\,dt)\) by smooth compactly supported \(p_n\) on \((-L,L)\), first excluding the finitely many zero points. This implies \(L^1\) convergence by Cauchy–Schwarz and \(\int w_\varepsilon^{-1}<\infty\). For a fixed smooth compactly supported \(\chi\) of integral one, replace \(p_n\) by \(r_n=p_n-(\int p_n)\chi\). Then \(\int r_n=0\) and \(r_n\to z'\) in weighted \(L^2\). The primitives \(z_n(t)=\int_{-L}^t r_n\) are smooth, vanish at both endpoints, and converge uniformly to \(z\). The functions \(z_n(x_1)b(x_2)\in H^1_0(U)\) converge to \(u^\varepsilon-g\) in weighted energy and \(L^1\). Approximation of each by \(C_c^\infty(U)\) in ordinary \(H^1\) also converges in weighted energy, since \(\mathbf a^\varepsilon\le I\). A diagonal sequence proves \(u^\varepsilon-g\in H^1_{\mathbf a^\varepsilon,0}(U)\). The usual energy identity proves uniqueness in that weighted class.\par
\subsection*{Failure of the asserted norm}
For \(|t|<\varepsilon/2\), [eq:volterra] (\#eq:volterra) gives \[v'(t)=c(\varepsilon/2)^\alpha |t|^{-\alpha}
       +O(|t|^{2-2\alpha}),\qquad c>0.\] Take \(s=1/10\) and nonzero nonnegative smooth compactly supported functions \(\psi\) on \((-1,1)\) and \(\eta\) on \((-L,L)\), with \(\int\eta b>0\). For \(0<\delta<\min\{\varepsilon/2,L\}\) let \(\phi_\delta(x)=\psi(x_1/\delta)\eta(x_2)\). Scaling the fractional seminorm yields \[\left\lVert\phi_\delta\right\rVert_{H^s(U)}\le C\delta^{1/2-s},\qquad
 \langle\partial_1u^\varepsilon,\phi_\delta\rangle
   =c_1\delta^{1-\alpha}+o(\delta^{1-\alpha}),\quad c_1>0.\] The constant-coefficient solution \(h\) is smooth on the compact support of these tests, so its pairing is \(O(\delta)\). Consequently \[\frac{|\langle\partial_1(u^\varepsilon-h),\phi_\delta\rangle|}
      {\left\lVert\phi_\delta\right\rVert_{H^s(U)}}
 \ge c_2\delta^{1/2+s-\alpha}\longrightarrow\infty,
 \qquad 1/2+s-\alpha=-1/5.\] These tests belong to \(C_c^\infty(U)\), so the conclusion applies to exactly the dual norm in (1.33), not merely to a boundary-sensitive variant. Volume normalization changes only a constant. Multiplication by the fixed invertible \(\overline{\mathbf s}^{1/2}\) cannot restore membership in \(H^{-s}\). Thus the first term in (1.20) is infinite and its right side is finite. The construction works for every fixed \(\varepsilon>0\), including arbitrarily small \(\varepsilon\). Increasing the finite homogenization scale cannot fix it.\par
\subsection*{Dependency, counterevidence, and repair boundary}
The source’s proof already exposes the restriction: p. 110 starts with \(s>\gamma/2\); (5.4), p. 86, requires a choice \(\gamma<\rho<2s\). The subsequent choices \(\rho=(1+s)/2\) and \(s=(1+\gamma)/2\) do not supply the stated full range. The printed sign in \((\rho-2s)^{-1}\) and the passage from decay of \(\mathcal E^2\) to decay of \(\mathcal E\) also require correction, but the counterexample does not depend on those textual mistakes.\par
This is a false principal conclusion for the reviewed version. It does not refute Theorem A, whose uniform ellipticity hypothesis excludes the example, or all other conclusions of Theorem B. A repair must restrict the exponent/topology or strengthen microscopic regularity. The written route suggests \(s>\gamma/2\) with \(\rho\in(\gamma,2s)\) and with the resulting dependence on the gap \(2s-\gamma\) tracked. That is a necessary proof-design boundary, not a certification of a replacement theorem with all original tail constants and quantifiers. Evidence: P01 and the complete saved periodic-counterexample proof, including the bounded helper’s coefficient-space and density cross-check.\par
\section*{A false large-contrast matrix bound}
Problem 2 (P02: Lemma 2.6, (2.81)–(2.82); Major).  The stated linear-in-contrast upper bounds are false under the explicit algebraic hypotheses of Lemma 2.6, pp. 35–36.\par
\subsection*{An exact admissible family}
Equation reference: eq:rankone. Let \(t>1\), \(d\ge2\), and take \[s_1=tI_d,\qquad s_{*,1}=I_d,\qquad k_1=\tfrac12(t-1)I_d.\] Tensor the following scalar \(2\)-by-\(2\) matrices with \(I_d\): \[E_1=\begin{pmatrix}(t+1)^2/4&-(t-1)/2\\-(t-1)/2&1\end{pmatrix},
 \quad
 E_{*,1}=\begin{pmatrix}(t+1)^2/(4t)&(t-1)/(2t)\\
                         (t-1)/(2t)&1/t\end{pmatrix}.\] They have precisely the form (2.78), are positive definite, and satisfy \[
 E_1-E_{*,1}
 =\frac{t-1}{4t}
   \begin{pmatrix}t+1\\-2\end{pmatrix}
   \begin{pmatrix}t+1&-2\end{pmatrix}\succeq0.\] The hypotheses allow arbitrary \(k_1\); skew-symmetry cannot be added silently. The contrast \(\widetilde\Theta_1\) in (2.80) is \(t\). The left side of (2.81) is \((t-1)^2\), whereas its claimed upper bound is \(2(t-1)\) for \(t>3\).\par
By [eq:rankone] (\#eq:rankone), one eigenvalue of \(E_{*,1}^{-1/2}E_1E_{*,1}^{-1/2}\) in each scalar block is \(1\). The determinant ratio is \(t^2\), so the other eigenvalue is \(t^2\). Its distance from the identity is therefore \(t^2-1\), not bounded by \(6(t-1)\) when \(t>5\). For \(t=10\), the two false inequalities read \[81\le18,\qquad 99\le54.\] All entries and determinant identities have also been checked by the saved reviewer-authored exact-rational script. No numerical spectral approximation is needed.\par
\subsection*{Failed step and a bounded repair}
On p. 36 the relevant congruence has top-left term \(4s_*^{-1/2}k\,s_*^{-1}k\,s_*^{-1/2}\) after removing the skew part of \(k\). The proof substitutes \(s_1^{-1}\) for the middle \(s_*^{-1}\) and invokes (2.85). Since \(s_1\succeq s_*\), that substitution reduces a positive term and does not give an upper bound for the original expression.\par
Equation reference: eq:quadraticrepair. For a general pair satisfying the lemma’s assumptions, put \(t=\left\lVert s_*^{-1/2}s_1s_*^{-1/2}\right\rVert\). The valid quadratic estimate (2.86) instead yields \[\left\lVert s_*^{-1/2}(k+k^t)s_*^{-1/2}\right\rVert^2\le(t-1)^2.\] It also gives the sharp replacement \[
 t-1\le
 \left\lVert E_*^{-1/2}EE_*^{-1/2}-I\right\rVert
 \le t^2-1.\] Indeed, after the same congruence, the positive semidefinite gap has diagonal-block norms bounded by \(t(t-1)\) and \(t-1\). For a positive semidefinite block matrix, block Cauchy–Schwarz bounds its norm by the sum of those two norms. This proves the upper bound; the family above attains it. The \(t=1\) case follows by continuity or directly from the positive block inequality. In particular, the printed linear bound still follows for \(1\le t\le5\). This is a verified algebraic repair, not a reproof of every downstream homogenization estimate.\par
\subsection*{Impact and counterevidence}
The source uses (2.82) in Lemma 2.7, p. 37, and invokes Lemma 2.6 in the comparison leading to (2.98), p. 39. These are quantitative inputs to the renormalization argument. The failure therefore affects a supporting estimate tied directly to the main proof framework. Correcting one inverse while retaining the lemma’s conclusions is insufficient.\par
The matrix family need not be realized as the coarse matrix of a stationary field: realization is not a hypothesis of Lemma 2.6. Conversely, it would be unjustified to claim that this family refutes every subsequent field-dependent conclusion. In particular, the intrinsic contrast \(\Theta\) used in (2.98) differs from \(\widetilde\Theta_1\): for this family it is \((t+1)^2/4\). The example does not refute a bound linear in that larger quantity. The remaining task is to audit each invocation with the correct contrast and regime. Evidence: P02, the exact matrix certificate, and source images of pp. 35–36.\par
\section*{The printed Poisson field is not elliptic}
Problem 3 (P03: (1.28), Proposition 1.1, and Appendix D.1; Major).  The additive counting formula in (1.28), p. 13, need not have positive coefficients, even at arbitrarily small positive inclusion intensity.\par
For ordinary convolution with a Poisson point measure, \(N_r(x)=(1_{B_r}*\omega_2)(x)\) counts points in \(B_r(x)\). Set \(d=2\), \(\rho_1=0\), \(\Lambda=1\), and \(\lambda=1/2\). The printed coefficient becomes \[\mathbf a(x)=\bigl(1-\tfrac12N_r(x)\bigr)I,\] where \(r=1\) in (1.28) and \(r=1/3\) in (D.1), p. 135. Let \(\rho_2>0\) be arbitrarily small. The event of exactly three points in \(B_{r/2}\) and no other points in \(B_{3r/2}\) has probability \[\exp(-\rho_2|B_{3r/2}|)\,
 \frac{(\rho_2|B_{r/2}|)^3}{6}>0.\] On that event, every \(x\in B_{r/2}\) sees exactly these three points within distance \(r\), apart from null boundaries. Consequently \(\mathbf a=-I/2\) on a set of positive measure. This violates both the ambient positivity requirement and the claimed uniform lower bound. A small-intensity restriction cannot remove an event of positive probability. Positive high-inclusion counts likewise need not obey the claimed upper bound \(\Lambda\).\par
Appendix D.1 explicitly changes only the radius, not counting to occupancy. Its use of \(\mathbf a\le\Lambda\) and of the inverse field on pp. 137–138 therefore cannot validate the printed model. A possible revised model replaces each count by its occupancy indicator \(1_{\{N_r>0\}}\). Its values are \(1,\Lambda,\lambda,\Lambda+\lambda-1\), all in \([\lambda,\Lambda]\). This repairs that elementary ellipticity failure, but changes the defined random field; the full stated coarse-ellipticity and tail estimates for the revised model have not all been rederived here. The objection concerns the claimed example, not the conditional abstract homogenization theorem. Evidence: P03 and the Poisson-event certificate; source pp. 13, 135, 137.\par
\section*{Matrix exponentiation does not provide the claimed ellipticity}
Problem 4 (P04: Proposition 1.3 and Proposition D.3; Major).  The unrestricted nonsymmetric matrix-exponential examples are invalid, and the uniform ellipticity estimate fails even within the positive symmetric-part class.\par
\subsection*{Failure of the Gaussian example}
On p. 15 the source explicitly allows a nonsymmetric real matrix-valued Gaussian field \(g\) and defines \(\mathbf a=\exp(hg)\). Let \[J=\begin{pmatrix}0&1\\-1&0\end{pmatrix},\qquad J^2=-I,
 \qquad g=F(I+J),\] where \(F\) is one nondegenerate mollified fractional Gaussian field with one of the stated Hurst parameters. Every matrix entry is \(F\) or \(-F\), with the required marginal law; the assumptions do not demand independent entries. Then \[\mathbf a=e^{hF}\bigl(\cos(hF)I+\sin(hF)J\bigr),\qquad
 \operatorname{sym}\mathbf a=e^{hF}\cos(hF)I.\] At a fixed point, \(F\) is a nondegenerate real Gaussian variable, so \(\cos(hF)<0\) has positive probability. By stationarity and Tonelli’s theorem, the expected measure of the negative set in any bounded region of positive volume equals that volume times the positive one-point probability. Thus there is positive probability of a positive-measure negative set. The failure cannot be removed by changing values on a null set; no continuity assumption is needed for this argument.\par
\subsection*{A counterexample that already belongs to \(\Omega\)}
One might try to read membership in \(\Omega\) as an implicit additional restriction. Proposition D.3, p. 145, is still false with its stated uniform parameters under that interpretation. For \(0<\theta<\pi/2\), take the deterministic constant field \[g_\theta=\theta J,\qquad h=\pi/2,\quad L=1,\qquad
 \mathbf a_\theta=e^{g_\theta}
 =\cos\theta\,I+\sin\theta\,J.\] The field is stationary and has finite range of dependence in the probabilistic sense, since its law is deterministic. Moreover \(\left\lVert g_\theta\right\rVert\le h\), so the weak \(\Gamma_2\) bound (D.26) holds for every \(\theta\) with the same \(h\). Now \(\operatorname{sym}\mathbf a_\theta=\cos\theta\,I>0\), and the constant field belongs to \(\Omega\).\par
The double-variable matrix \(A\) in (2.8), p. 20, has lower-right block \(\sec\theta\,I\). The bounds (2.22), p. 22, squeeze the coarse matrix of a constant field to that same pointwise \(A\) on every domain. Taking \(n=m\) in (P2\(\dagger\)), (2.71), p. 33, shows that any finite minimal scale must eventually give \(A(\square_m)\preceq E_0\). But Proposition D.3 claims \(E_0=C\exp(18h^2)I\), with \(C=C(d)\) independent of \(\theta\). As \(\theta\uparrow\pi/2\), \(\sec\theta\) is unbounded. No finite universal \(C\) works, and delaying the scale cannot improve a constant matrix. The asserted tail bound forces the scale to be finite, so it does not remove this contradiction.\par
The specific estimate used near p. 149, \(\left\lVert A\right\rVert\le(1+\left\lVert\mathbf a\right\rVert^2)\left\lVert\mathbf a^{-1}\right\rVert\), also fails: the right side is \(2\) for this orthogonal rotation, whereas \(\left\lVert A\right\rVert\ge\sec\theta\).\par
\subsection*{Impact and repair boundary}
Invertibility and bounded singular values of \(\exp(g)\) do not control its symmetric part. The negative example addresses the stated class; the positive rotation family addresses the quantitative claim even after positivity is imposed. Both are relevant counterevidence to a mere implicit-ellipticity defense. Symmetric \(g\) would make \(\exp(g)\) positive definite, with inverse bounds controlled by \(e^{\left\lVert g\right\rVert}\); a sufficiently quantitative sector condition is another possible starting point. Neither observation certifies the source’s full probabilistic statement, constants, or mixed-entry dependence under such a revision. Evidence: P04, the rotation certificate, and pp. 15, 145, 149.\par
\section*{A false tail bound for independent sums}
Problem 5 (P05: Lemma C.2, (C.15); Major).  For \(1<\sigma<2\), the large-deviation term in (C.15), p. 129, omits essential dependence on the number of summands.\par
Let \(\sigma=3/2\), \(m=100\), and \(t=10^8\). Take independent \(X_i=\epsilon_iV_i\), where \(\epsilon_i\) are independent Rademacher signs and \(\mathbb P(V_i>u)=\exp(-u^{3/2})\) for \(u\ge0\). Then \(\mathbb EX_i=0\) and \(\mathbb P(|X_i|>u)=\exp(-u^{3/2})\), so every required weak Orlicz bound holds with scale one.\par
The event \(X_i>t/m\) for all \(i\) implies \(\sum_iX_i>t\) and has logarithm \[-100\log2-10^{11}>-100000000100.\] The two quantities inside the maximum on the right of (C.15) have logarithms \[-\frac{t^2}{40m}=-2500000000000,\qquad
 \log1024-\frac{t^{3/2}}3
 <10-\frac{10^{12}}3.\] Both are strictly smaller than the logarithm of the exhibited event. Thus the claimed upper bound is false. These comparisons use exact integer bounds and \(\log2<1\); no floating-point evaluation of exponentially small probabilities is involved.\par
The source starts with the correct Chernoff product (C.17), but on p. 131 inserts a one-variable moment-generating-function bound without retaining its \(m\)th power. The same all-positive event, at general \(m\), has leading logarithm \(-m^{1-\sigma}t^\sigma\). A uniform large-deviation exponent cannot be independent of \(m\) in the claimed way.\par
This is a substantive error in an explicit probabilistic ingredient of the quantitative framework. It does not by itself refute the weaker normalized Orlicz estimate (C.16), all of Appendix C, or Theorem A. Those conclusions would require a valid independent-sum argument, including its actual constants. A corrected Chernoff calculation must retain the product and optimize at the appropriate \(t/m\) scale. A full replacement with the printed numerical constants has not been verified here. Evidence: P05 and the exact tail certificate, source pp. 129–131.\par
\section*{The renormalization reduction does not meet its stated premises}
Problem 6 (P06: scale selection in Section 3.1; Major proof dependency).  The determinant argument and its application on p. 54 have incompatible normalizations and scale budgets. A further strengthened hypothesis is used on pp. 55–57 without following from the stated one.\par
\subsection*{Determinant normalization and the missing scale factor}
The source defines \(\Xi_j=\det(\overline A_j)^{1/d}\) in (2.107), p. 41. If \(B\succeq C>0\) but \(B\npreceq(1+\delta_1)C\), at least one eigenvalue of \(C^{-1/2}BC^{-1/2}\) exceeds \(1+\delta_1\), and all are at least one. The determinant ratio therefore gives only \[\bigl(\det C/\det B\bigr)^{1/d}
 \le(1+\delta_1)^{-1/d},\] whereas Lemma 3.4 uses \((1+\delta_1)^{-1}\). The discrepancy is witnessed exactly in the same block-algebraic class. Take \(d=2\), \(\delta_1=\sigma=1/2\), \(h=1\), \(N=3\), \(r=25/16\), and \[\overline A_j=\operatorname{diag}(r^{3-j},1,1,1),\qquad 0\le j\le3.\] These matrices are decreasing and have the prescribed block form with \(k=0\) and \(s\succeq s_*=I\). The printed count \(N\ge\lceil4\log2\rceil=3\) holds. Every required consecutive Loewner-closeness alternative fails, since \(r>3/2\), but \(\Xi_3/\Xi_0=r^{-3/2}=64/125>1/2\). This refutes the matrix inference used by the proof. It is not a claim that this finite sequence has been realized by an admissible random coefficient field.\par
A sufficient matrix-only repair is \[N\ge h\left\lceil2d\,\delta_1^{-1}|\log\sigma|\right\rceil,\] using \(\log(1+\delta_1)\ge\delta_1/2\). The application (3.9) has a second problem: it sets \(\delta_1=\delta\sigma^2/2\) and \(h=2n\), but takes \[N/h=\left\lceil2\delta^{-1}|\log\sigma|\right\rceil.\] Even the printed lemma requires \(\lceil4\delta^{-1}\sigma^{-2}|\log\sigma|\rceil\); the repaired matrix argument requires \(\lceil4d\delta^{-1}\sigma^{-2}|\log\sigma|\rceil\). For \(d=2\) and \(\delta=\sigma=1/2\), the three counts are respectively \(3\), \(23\), and \(45\). The missing \(\sigma^{-2}\) cannot be absorbed into a universal constant as \(\sigma\downarrow0\).\par
\subsection*{The strengthened hypothesis is not automatic}
The first term of (3.5), p. 54, is \(\Pi3^{-(1-\beta)\ell/4}\). The purported repeat (3.11), p. 55, instead uses \[\frac{K_{\Psi_{\mathcal S}}^{3\gamma}}{1-\gamma}
 \Pi3^{-(1-\beta)\ell/4},\] and that stronger estimate is used on p. 57. Here is a complete parameter check that survives arbitrarily small \(\delta\). Let \(q\ge1\) be any integer and set \[\begin{gathered}
 d=2,\quad\beta=7/8,\quad\gamma=1/2,\quad\nu=1,\quad
 \Pi=2,\quad K_\Psi=3,\quad K_{\Psi_{\mathcal S}}=3^{100},\\
 \delta=3^{-q},\qquad \sigma=1/2,\qquad
 \ell=3296+32q,\qquad m=100\ell.
 \end{gathered}\] The four terms in (3.5) are \[2\,3^{-103-q},\quad
 2\,3^{-95-q},\quad
 256\,3^{-96-32q},\quad
 2\,3^{-1648-16q}.\] Each is at most \(\delta\sigma^2=\tfrac14\,3^{-q}\). The strengthened first term in (3.11) is instead \(4\,3^{47-q}\), whose ratio to that threshold is \(16\,3^{47}>1\).\par
These parameters are consistent with an actual law: take the constant field \(\mathbf a=\operatorname{diag}(2,1)\), a Dirac measure, \(E_0=\operatorname{diag}(2,1,1/2,1)\), \(\mathcal S=0\), and \(\Psi=\Psi_{\mathcal S}=\exp(t)\). The allowed growth constants need not be minimal, so \(K_{\Psi_{\mathcal S}}=3^{100}\) is admissible. The contrast is \(\Theta=1\), \(\Pi=2\), and (3.6) holds exactly. For any positive \(\delta_0(d)\), choose \(q\) large enough. The strengthened premise still fails. The proposition’s conclusion is true for this constant field; the example proves a failed implication between proof premises, not a false homogenization conclusion.\par
\subsection*{Impact, a surviving iteration, and repair limits}
These steps are the reduction to the central one-step estimate. The direct block-count repair introduces a \(\sigma^{-2}|\log\sigma|\) cost. It does not by itself recover the displayed uniform-in-\(\sigma\) scale in Proposition 3.2. The strengthening of (3.5) also needs a corresponding scale choice.\par
There is useful counterevidence to an overbroad negative conclusion. At a fixed target \(\sigma_0(d)>0\), the extra factors are dimensional constants. Conditional on the other one-step hypotheses, the logarithmic-in-contrast step cost and logarithmic iteration count still have the order needed for an \(\exp(C\log^2\text{contrast})\) scale.\par
The shifted-product iteration on pp. 72–73 can also be made rigorous in its needed range by a positive part. If decreasing \(x_j,y_j\ge1\) satisfy \[\min\{(x_{j+1}-1)/x_j,\ y_{j+1}/y_j\}\le\sigma,\qquad
 0<\sigma\le1/2,\] set \(u_j=(x_j-1-2\sigma)_+\). Either the first alternative gives \(u_{j+1}\le\sigma u_j\), or the second gives \(u_{j+1}y_{j+1}\le\sigma u_jy_j\) by monotonicity. Thus \(u_ky_k\le\sigma^k x_0y_0\le\sigma^k\Theta^2\). For \(k\ge\lceil2\log_2\Theta\rceil+1\), this yields \(x_k\le1+3\sigma\). The unnecessary untruncated product estimate is not being accepted as a proof outside that stopping regime.\par
This is a material proof-closure concern, with exact local repairs and a precise limit on the preserved asymptotic claim. No proof is given here of a false Theorem A or of the full arbitrary-\(\sigma\) replacement. Evidence: P06, the two exact arithmetic certificates, and source pp. 41, 54–55, 57, 72–73.\par
\section*{The stated finite-range decomposition is not constructed}
Problem 7 (P07: Appendix D.2.2, (D.18)–(D.22); Major proof dependency).  The printed cutoffs cannot exist, and the subsequent heat-convolution definition does not produce the finite-range fields required by the example-verification argument.\par
On pp. 142–143 the annuli \(B_{3^n}\setminus B_{3^{n-1}}\), \(n\in\mathbb Z\), tile punctured space apart from boundaries. The requirements \[1_{B_{3^n}\setminus B_{3^{n-1}}}\le\eta_n,\qquad
 \eta_n\ge0,\qquad \sum_n\eta_n=1\] force \(\eta_n=1\) inside its own annulus and zero inside every other one. Approaching an annular boundary from the two sides gives a jump. No continuous, and hence no smooth compactly supported, family satisfies the stated conditions.\par
There is a separate analytic obstruction. Write \(\alpha=d/4-\sigma/2\in(0,d/4)\) and let \[k_\alpha(x)=
 \frac{\Gamma(d/2-\alpha)}
 {4^\alpha\pi^{d/2}\Gamma(\alpha)}
 |x|^{2\alpha-d}\] be the heat-potential kernel from (D.13). The definition immediately before (D.22) integrates a heat convolution with \(\eta_n\); its white-noise kernel is \(k_\alpha*\eta_n\), not the pointwise product \(k_\alpha\eta_n\) written in (D.22). For a nonzero nonnegative smooth annular cutoff, \(k_\alpha*\eta_n\) is strictly positive everywhere and lies in \(L^2\). Its covariance with every translate is strictly positive. The resulting stationary Gaussian field has no finite dependence range. The product kernel, in contrast, is compactly supported. Thus the printed identity does not establish the required property.\par
Even for the intended product kernels, use of the same white noise does not make entire fields at different scales independent. Translate one nonzero kernel so that it overlaps a positive open part of another. Their cross-covariance is then positive. Disjoint centered annuli could at most supply a fixed-center observation. Cross-scale independence is fortunately not an assumption of Propositions D.2 and D.4, so it need not be restored.\par
\subsection*{A valid local construction and its boundary}
Choose a smooth radial nonincreasing \(\chi\) equal to one on \(B_1\) and zero outside \(B_{3/2}\), and set \[\eta_n(x)=\chi(3^{-n}x)-\chi(3^{1-n}x),\qquad
 \widetilde F_{\sigma,n}(x)
 =\int k_\alpha(y-x)\eta_n(y-x)\,W(dy).\] The cutoffs are nonnegative and smooth, sum to one off the origin, and have the required dilation relation; the impossible indicator lower bound is removed. Each product kernel is supported in \(B_{(3/2)3^n}\), so the field has dependence range at most \(3^{n+1}\). White-noise scaling gives \[\widetilde F_{\sigma,n}\ \stackrel{\rm law}{=}\
 3^{-n\sigma}\widetilde F_{\sigma,0}(3^{-n}\,\cdot).\] The field and gradient variance scales are respectively \(3^{-2n\sigma}\) and \(3^{-2n(1+\sigma)}\).\par
For a Schwartz test function, the convolved partial kernels converge in \(L^2\) to \(k_\alpha*\psi\): the singularity is locally integrable because \(\sigma<d/2\), and the decay at infinity is square-integrable because \(\sigma>0\). Dominated convergence and white-noise isometry give the distributional decomposition. No cross-scale independence is asserted.\par
This proves a qualitative replacement with finite range, scaling, and distributional convergence. It does not certify all uniform dependence on \(\sigma\) in the subsequent supremum estimates, nor the mixing assertions referred to external work. The original construction fails to supply a necessary premise for the Gaussian-field application, which is the reason for its substantive classification. The finding is not a counterexample to Proposition 1.2 itself. Evidence: P07, source pp. 142–143, and the saved reconstruction.\par
\section*{Coverage, positive checks, and unresolved dependencies}
The complete text of all 152 source pages was read across the root conversation and the three bounded native helpers. That coverage is distinct from proving every assertion anew. The ledger identifies the separate ranges, actually viewed source images, exact arithmetic, and proof reconstructions.\par
Besides the decisive checks above, the following bounded arguments were followed or reconstructed: the coarse matrix variational sandwich and constant-field case; the invariance needed for the matrix congruences; the valid quadratic small-contrast bound; the weighted solution-space closure for Problem 1; the positive-part high-contrast iteration; and the local truncation argument of Proposition C.3. For the last item, the intermediate integral on p. 132 is \(1+\lambda\le2\), and the moment-generating-function coefficient in (C.23) still follows from \((1+e+2M+2C_\Psi)/2\le2+M+C_\Psi\). A conservative corrected version of (C.29) is proved in the companion check, rather than treating its coefficient slips as another substantive problem.\par
The following remain explicit limits. The complete parameter cost of repairing the high-contrast estimates has not been propagated. The proof of Proposition 4.2 uses a small-input-contrast premise through Lemma 4.5 that is not identical to small coarse contrast at scale zero; a sufficient additional renormalization has not been fully quantified here. The large-\(k\) matrix estimate in Lemma 4.4 and the \(D=0\) endpoint of Theorem 6.1 remain questions requiring evidence, not confirmed false-theorem claims. The external results attributed to [AK24], [JK95], [DLV18], and [FP97] in the source were identified, but their complete primary proofs were not independently rederived or represented as fully verified. These limits do not enter any of the explicit counterexamples above as unproved assumptions.\par
The primary-source lookup confirmed the v3 record and the publisher identity. A targeted search of the arXiv, publisher, and author domains did not locate an erratum or correction for the identified version. This limited search is not proof that no later correction exists. No author adoption, author response, or human editorial decision has been obtained.\par
\section*{Conclusion}
The reviewed author manuscript requires Major revision on the evidence given here. Problems 1–5 include false stated conclusions under their printed hypotheses; Problems 6–7 establish invalid proof dependencies with their narrower consequences stated explicitly. The precise uniformly elliptic main theorem is not refuted by these examples. Its full proof should be assessed after the supporting estimates and parameter reductions have been repaired.\par
Dimension | Recorded outcome\par
Substantive findings | Seven Major concerns, with separate proofs and scope limits\par
Verified repairs | Matrix bound \(t^2-1\); periodic-example density; corrected local finite-range representation; bounded iteration and tail corrections\par
Not certified | A replacement for all main theorems with the original constants, quantifiers, and tail parameters\par
Author adoption | Unknown; no contact or adoption asserted\par
Scientific adjudication | Pending\par
Human reviews | None recorded\par
Independent acceptance | Not performed in this initial review\par
\section*{Supporting checks}
Each P01–P07 directory contains a hash-bound evidence manifest linking the reviewed PDF, the relevant source-page images, the analytic proof, and any actual reviewer-authored calculation. The complete calculations use exact rational or integer comparisons; they do not substitute finite numerical samples for universal arguments.\par
> `81aedeb87003afb4813914f60dd54f88e688bce01529256418974e35b616b5b6`.\par
Actual compilation, dependency copies, compiler and input hashes, stdout/stderr, retained failed attempts, report-page rendering, and actual page-inspection observations are recorded in the companion provenance. The existing authorized Tectonic runtime is used in untrusted mode with shell escape disabled. Only reviewer-written checks were executed. The source manuscript was never compiled or run as author code. Earlier drafts and real failures remain preserved.\par
The root scientific conversation integrated the proofs, wrote and revised the English report, ran the actual report build, and carried out the report-page inspection. Three native helpers in the same conversation supplied bounded foundation, renormalization, and appendix checks. They are not separate acceptance reviewers.\par
The account/mode/model labels were supplied by the external controller’s UI observation as Personal account, Chat, GPT-6, Pro. No backend model identity was independently measured or inferred from either label. The earlier controller-contract helper supplied only an operational role audit and is not scientific evidence.\par
Local formula slips with directly demonstrated corrections, including the scale normalization in Lemma 2.2, a matrix symbol in (2.27), and the exponent sign in (5.74), are retained in the bounded check records. They are excluded from the substantive issue count. No claim is made that the seven listed concerns exhaust all possible errors in this long manuscript.\par
9 S. Armstrong and T. Kuusi, Renormalization group and elliptic homogenization in high contrast, arXiv:2405.10732v3, September 6, 2025, 152 pp. https://arxiv.org/abs/2405.10732v3. Springer Nature, publisher record for the same title, Inventiones mathematicae 242 (2025), 895–1086, DOI https://doi.org/10.1007/s00222-025-01370-9. The record was consulted as a subscription preview; the journal PDF was not obtained.\par
\end{document}
